%%%%%%%%%%%%%%%%%%%%%%% file example.tex %%%%%%%%%%%%%%%%%%%%%%%%%%
%
% This is a template file for Journal of Physics: Conference Series
%
% Copy it to a new file with a new name and use it as the basis
% for your article
%
%%%%%%%%%%%%%%%%%%%%%%%%%% IOP Publishing %%%%%%%%%%%%%%%%%%%%%%%%%
%

%% The first command in your LaTeX source must be the \documentclass command.
%%
\documentclass[a4paper]{jpconf}

%
% Put here some packages required or/and some personnal commands
%
\usepackage{amsmath}
\usepackage{hyperref}
\hypersetup{hidelinks}
\usepackage{graphicx}
\usepackage{enumitem}

\usepackage{iopams}
\usepackage{cite}
\bibliographystyle{iopart-num}   % IOP citation style
\newcommand{\ees}{\textit{IOP~Conference Series: Earth and Environmental Science}}
%
%% end of the preamble, start of the body of the document source.
\begin{document}
%
%% The "title" command
\title{A better way to format your ICon-MaSTEd paper for \ees{}}
%
% subtitle is optional
%
%%%\subtitle{Do you have a subtitle?\\ If so, write it here}

%% The "author" command and its associated commands are used to define
%% the authors and their affiliations.

%%
%% The "author" command and its associated commands are used to define
%% the authors and their affiliations.
\author{S O Semerikov$^{1,2,3}$, S M Chukharev$^4$, S I Sakhno$^2$, A M Striuk$^{2,5}$, Andrii V Iatsyshin$^{6,7}$, S V Klimov$^4$, V V Osadchyi$^8$, T A Vakaliuk$^{9,3,1}$, P P Nechypurenko$^1$, O V Bondarenko$^1$ and H B Danylchuk$^{10}$}

\address{$^{1}$ Kryvyi Rih State Pedagogical University, 54 Gagarin Ave., Kryvyi Rih, 50086, Ukraine}
\address{$^{2}$ Kryvyi Rih National University, 11 Vitalii Matusevych Str., Kryvyi Rih, 50027, Ukraine}
\address{$^{3}$ Institute for Digitalisation of Education of the NAES of Ukraine, 9 M.~Berlynskoho Str., Kyiv, 04060, Ukraine}
\address{$^{4}$ National University of Water and Environmental Engineering, 11 Soborna Str., Rivne, 33028, Ukraine}
\address{$^{5}$ Academy of Cognitive and Natural Sciences, 54 Gagarin Ave., Kryvyi Rih, 50086, Ukraine}
\address{$^{6}$ State Institution ``The Institute of Environmental Geochemistry of National Academy of Sciences of Ukraine'', 34a Palladin Ave., Kyiv, 03142, Ukraine}
\address{$^{7}$ G.E. Pukhov Institute for Modelling in Energy Engineering of NAS of Ukraine, 15 General Naumova Str., Kyiv, 03164, Ukraine}
\address{$^{8}$ Bogdan Khmelnitsky Melitopol State Pedagogical University, 20 Hetmanska Str., Melitopol, 72300, Ukraine}
\address{$^{9}$ Zhytomyr Polytechnic State University, 103 Chudnivska Str., Zhytomyr, 10005, Ukraine}
\address{$^{10}$ Bohdan Khmelnytsky National University of Cherkasy, 81 Shevchenko Blvd., Cherkasy, 18031, Ukraine}

\ead{semerikov@gmail.com, konf.knu@gmail.com , budfac@gmail.com, andrey.n.stryuk@gmail.com, iatsyshyn.andriy@gmail.com, s.v.klimov@nuwm.edu.ua, poliform55@gmail.com, tetianavakaliuk@gmail.com, acinonyxleo@gmail.com, bondarenko.olga@kdpu.edu.ua, abdanilchuk@gmail.com}


%% The abstract is a short summary of the work to be presented in the
%% article.
\begin{abstract}
A clear and well-documented \LaTeX{} document is presented as an
  article formatted for publication by \iopp in a \ees{}. Based on the ``jpconf'' document class, this article presents and explains many of the common variations, as well as many of the formatting elements an author may use in the preparation of the documentation of their work. An article abstract should not normally exceed 200 words in a single paragraph. This template don't use any keywords.
\end{abstract}
%

\section{On the \TeX{}}

Already the 2nd proceedings of the ICSF conference is being prepared for publication in the \ees{}. Working on past proceedings led the editorial board closer to the idea that it is easier to reject a poorly formatted article rather than spend time editing it and delaying the publication of the complete proceedings volume. As a result, authors are highly recommended to read this handbook from beginning to end before beginning work on their piece. Failure to follow the formatting guidelines will result in the article being rejected at the review stage or even sooner.

Traditionally, we use \LaTeX{} templates for the conference proceedings due to a lot of reasons, the main of them is the attempt to decrease an extra amount of editing efforts for the proceedings editors. You can freely use any \LaTeX{} compatible typesetting system (e.g., TeXstudio + TeX Live is a good choice for any operating systems), but if you don't to be involved into the \LaTeX{} system administration, we propose to use a cloud based \LaTeX{} editors like Overleaf.
After registering at \url{https://www.overleaf.com}, you can start your paper revision with this template using `New Project' -- `Upload Project' menu (\fref{overleaf_1}).

\begin{figure}[h]
\centering
\includegraphics[width=0.75\linewidth]{img/overleaf_1}
\caption{How to upload your project to Overleaf, part 1.}
\label{overleaf_1}       % Give a unique label
\end{figure}

The next step is to select the template archive -- you can download it from the conference website (\fref{overleaf_2}, \fref{overleaf_3}).

\begin{figure}[h]
\centering
\includegraphics[width=0.75\linewidth]{img/overleaf_2}
\caption{How to upload your project to Overleaf, part 2.}
\label{overleaf_2}       % Give a unique label
\end{figure}

\begin{figure}[h]
\centering
\includegraphics[width=\linewidth]{img/overleaf_3}
\caption{Overleaf, online \LaTeX{} editor.}
\label{overleaf_3}       % Give a unique label
\end{figure}

To get a camera-ready version of your paper in PDF, you can click to 'Download PDF' icon or use 'Menu' to download both LaTeX source files (ZIP) and camera-ready version (PDF) (\fref{overleaf_4}).

\begin{figure}[h]
\centering
\includegraphics[width=0.75\linewidth]{img/overleaf_4}
\caption{How to download your project from Overleaf.}
\label{overleaf_4}       % Give a unique label
\end{figure}


The most-often recommended tutorials are ``The (Not So) Short Introduction to \LaTeX{}2$\varepsilon$''. (\url{https://www.ctan.org/tex-archive/info/lshort/}) and ``Learn LaTeX in 30 minutes''  (\url{https://www.overleaf.com/learn/latex/Learn_LaTeX_in_30_minutes}).

\section{On the template}
\label{template_about}

\verb"jpconf" requires \LaTeXe\ and  can be used with other package files such
as those loading the AMS extension fonts 
\verb"msam" and \verb"msbm" (these fonts provide the 
blackboard bold alphabet and various extra maths symbols as well as 
symbols useful in figure captions); an extra style file \verb"iopams.sty" is 
provided to load these packages and provide extra definitions for bold Greek letters. 

The \verb"iopart-num" \BibTeX{} style is intended for use in
preparing manuscripts for Institute of Physics Publishing journals,
including \ees{}.  It provides numeric citation with
Harvard-like formatting.

If you are new to publishing with \ees{}, this document is a valuable
guide to the process of preparing your work for publication.

To begin the use of template, you need to:

\begin{enumerate}
 \item Download and unpack \LaTeX{} template  \url{http://cms.iopscience.iop.org/alfresco/d/d/workspace/SpacesStore/a83f1ab6-cd8f-11e0-be51-5d01ae4695ed/LaTeXTemplates.zip}
 \item Download and unpack \BibTeX{} style  \url{https://mirrors.ctan.org/biblio/bibtex/contrib/iopart-num.zip}
 \item Copy \verb|jpconf.cls| and \verb|jpconf11.clo| from unpacked \verb|LaTeXTemplates.zip|, and \verb|iopams.sty| and \verb|iopart-num.bst| from unpacked \verb|iopart-num.zip| to a place where \LaTeX{} can find them or simply copy them in the same directory as the source file of the article.

\end{enumerate}


\section{Modifications}
\label{nomodify}

Modifying the template -- including but not limited to: adjusting
margins, typeface sizes, line spacing, paragraph and list definitions,
and the use of the \verb|\vspace| command to manually adjust the
vertical spacing between elements of your work -- is not allowed.

\subsection{Headers, footers and page numbers}
Authors should {\it not} add headers, footers or page numbers to the pages of their article -- they will be added by \iopp\ as part of the production process.


\section{Front matter}
\label{fpage}

\subsection{Title information}

The titles of papers should should all use the regular English style: the first letter of the title should be capitalized with the rest in lower case. Use the \verb|title| command to define the title of your work. Do not
insert line breaks in your title.
\begin{verbatim}
\title{A better way to format your document for \ees{}}
\end{verbatim}


\subsection{Authors and affiliations}

The next information required is the list of all authors' names followed by 
the affiliations. For the authors' names type \verb"\author{#1}", 
where \verb"#1" is the 
list of all authors' names. The style for the names is initials then
surname, with a comma after all but the last 
two names, which are separated by `and'. Initials should {\it not} have 
full stops. 

The addresses of the authors' affiliations follow the list of authors. 
Each address should be set by using
\verb"\address{#1}" with the address as the single parameter in braces. 
If there is more 
than one address then a superscripted number, followed by a space, should come at the start of
each address. In this case each author should also have a superscripted number or numbers following their name to indicate which address is the appropriate one for them.

Please ensure that affiliations are as full and complete as possible and include the department, institution, full postal address and postal index, and country. 
If the authors are at different addresses, numbered superscripts should be used after each surname to reference an author to his/her address.
Multiple authors may share one affiliation. 

Please also provide e-mail addresses for any or all of the authors using an \verb"\ead{#1}" command after the last address. \verb"\ead{#1}" provides the text Email: so \verb"#1" is just the e-mail address or a list of emails.  

\begin{verbatim}
\author{S O Semerikov$^{1,2,3}$, V N Soloviev$^{1}$ and A E Kiv$^{4}$}

\address{
    $^{1}$ Kryvyi Rih State Pedagogical University, 
    54 Gagarin Ave., Kryvyi Rih, 50086, Ukraine
}
\address{
    $^{2}$ Kryvyi National University, 
    11 Vitalii Matusevych Str., Kryvyi Rih, 50027, Ukraine
}
\address{
    $^{3}$ Institute for Digitalisation of Education of the NAES of Ukraine, 
    9 M. Berlynskoho Str., Kyiv, 04060, Ukraine
}
\address{
    $^{4}$ Ben-Gurion University of the Negev, 
    P.O.B. 653, Beer Sheva, 8410501, Israel
}

\ead{
    semerikov@gmail.com,
    vnsoloviev2016@gmail.com,
    kiv.arnold20@gmail.com
}
\end{verbatim}

\subsection{Abstract and keywords}

The abstract follows the addresses and
should give readers concise information about the content 
of the article and should not normally exceed 200 
words. {\bf All articles must include an abstract}. To indicate the start 
of the abstract type \verb"\begin{abstract}" followed by the text of the 
abstract.  The abstract should normally be restricted 
to a single paragraph and is terminated by the command
\verb"\end{abstract}"

\begin{verbatim}
\begin{abstract}
  This is an abstract.
\end{abstract}
\end{verbatim}

Do not enter key words for this journal.

The command \verb"\maketitle" is not
required.

\subsection{Sections and subsections}\label{subsection}

\begin{verbatim}
\section{This is the section title}
\subsection{This is the subsection title}\label{subsection}
\end{verbatim}

Cross references to other sections in the text should, where possible, 
be made using 
labels (see \sref{subsection}) 
but can also be made manually.

\begin{verbatim}
\subsection{Sections and subsections \label{subsection}}
\end{verbatim}

\section{Tables}

Tables should be numbered sequentially throughout the text and referred to in the text 
by number (table 1, etc, {\bf rather than} tab. 1). 
Each table should be a float and be positioned within the text at the most convenient place 
near to where it is first mentioned in the text. It should have an 
explanatory caption which should be as concise as possible. 
Captions should be placed at the top of the table and should have a full stop (period) at the end. 

\subsection{The basic table format}
The standard form for a table is:
\begin{verbatim}
\begin{table}
\caption{Table caption.}
\label{label}
\begin{center}
\begin{tabular}{llll}
\br
Head 1&Head 2&Head 3&Head 4\\
\mr
1.1&1.2&1.3&1.4\\
2.1&2.2&2.3&2.4\\
\br
\end{tabular}
\end{center}
\end{table}
\end{verbatim}

The above code produces \tref{ex}.

\begin{table}[h]
\caption{Table caption.}
\label{ex}
\begin{center}
\begin{tabular}{llll}
\br
Head 1&Head 2&Head 3&Head 4\\
\mr
1.1&1.2&1.3&1.4\\
2.1&2.2&2.3&2.4\\
\br
\end{tabular}
\end{center}
\end{table}

Points to note are:
\medskip
\begin{enumerate}[label=\arabic*.]
\item The caption comes before the table.

\item The normal style is for tables to be centred in the same way as equations. This is accomplished by using \verb"\begin{center}" \dots\ \verb"\end{center}".

\item The default alignment of columns should be aligned left.

\item Tables should have only horizontal rules and no vertical ones. The rules at the top and bottom are thicker than internal rules and are set with \verb"\br" (bold rule). The rule separating the headings from the entries is set with
\verb"\mr" (medium rule). These commands do not need a following double backslash.

\item Numbers in columns should be aligned as appropriate, usually on the decimal point;
to help do this a control sequence \verb"\lineup" has been defined 
which sets \verb"\0" equal to a space the size of a digit, \verb"\m"
to be a space the width of a minus sign, and \verb"\-" to be a left
overlapping minus sign. \verb"\-" is for use in text mode while the other
two commands may be used in maths or text.
(\verb"\lineup" should only be used within a table
environment after the caption so that \verb"\-" has its normal meaning
elsewhere.) See \tref{tabone} for an example of a table where
\verb"\lineup" has been used.
\end{enumerate}

\begin{table}[h]
\caption{\label{tabone}A simple example produced using the standard table commands 
and $\backslash${\tt lineup} to assist in aligning columns on the 
decimal point. The width of the 
table and rules is set automatically by the 
preamble.} 

\begin{center}
\lineup
\begin{tabular}{*{7}{l}}
\br                              
$\0\0A$&$B$&$C$&\m$D$&\m$E$&$F$&$\0G$\\
\mr
\0\023.5&60  &0.53&$-20.2$&$-0.22$ &\01.7&\014.5\\
\0\039.7&\-60&0.74&$-51.9$&$-0.208$&47.2 &146\\
\0123.7 &\00 &0.75&$-57.2$&\m--   &--  &--\\ 
3241.56 &60  &0.60&$-48.1$&$-0.29$ &41   &\015\\ 
\br
\end{tabular}
\end{center}
\end{table}
 
You can find a lot of examples at \href{https://www.overleaf.com/learn/latex/Tables}{\textit{Overleaf documentation on tables}}. 

\section{Math equations}

You may want to display math equations in three distinct styles:
inline, numbered or non-numbered display.  Each of the three are
discussed in the next sections.

Equations may be numbered sequentially throughout the text (i.e., (1), (2), (3), ...) 
or numbered by section (i.e., (1.1), (1.2), (2.1), ...) depending on the author’s personal preference. 
In articles with several appendices equation numbering by section is useful in the appendices even 
when sequential numbering has been used throughout the main body of the text: for example, 
A.1, A.2 and so forth. When referring to an equation in the text, always put the equation number 
in brackets -- e.g. ‘as in equation (2)’ or ‘as in equation (2.1)’ -- and always spell out the word ‘equation’ 
in full, e.g. ‘if equation (5) is factorized’; do not use abbreviations such as ‘eqn.’ or ‘eq.’.

\subsection{Inline (in-text) equations}

A formula that appears in the running text is called an inline or
in-text formula.  It is produced by the \verb|math| environment,
which can be invoked with the usual
\verb|\begin| \ldots \verb|\end| construction or with
the short form \verb|$| \ldots \verb|$|. You can use any of the symbols
and structures, from $\alpha$ to $\omega$;
this section will simply show a few
examples of in-text equations in context. Notice how this equation:
\begin{math}
  \lim\limits_{n\rightarrow \infty} \frac{1}{n} = 0,
\end{math}
set here in in-line math style, looks slightly different when
set in display style.  (See next subsection).

\subsection{Display equations}

A numbered display equation -- one set off by vertical space from the
text and centered horizontally -- is produced by the \verb|equation|
environment. An unnumbered display equation is produced by the
\verb|displaymath| environment (or \verb|equation*| with \verb|amsmath| package).

Again, in either environment, you can use any of the symbols and
structures available in \LaTeX{}; this section will just give a couple
of examples of display equations in context.  First, consider the
equation, shown as an inline equation above:
\begin{verbatim}
\begin{equation}
  \lim_{n\rightarrow \infty} \frac{1}{n} = 0.
\end{equation}
\end{verbatim}
\begin{equation}
  \lim_{n\rightarrow \infty} \frac{1}{n} = 0.
\end{equation}
Notice how it is formatted somewhat differently in
the \verb|displaymath|
environment.  Now, we'll enter an unnumbered equation:
\begin{verbatim}
\begin{displaymath}
  S_{n} = \sum_{i=1}^{n} x_{i} ,
\end{displaymath}
\end{verbatim}
\begin{displaymath}
  S_{n} = \sum_{i=1}^{n} x_{i} ,
\end{displaymath}
and follow it with another numbered equation:
\begin{verbatim}
\begin{equation}\label{lim}
  \lim_{x \to 0} (1 + x)^{1/x} = e
\end{equation}
\end{verbatim}
\begin{equation}\label{lim}
\lim_{x \to 0} (1 + x)^{1/x} = e
\end{equation}
just to demonstrate \LaTeX's able handling of numbering.

Usually, equations should be centred and should be numbered with the number on the right-hand side. (You can find an additional examples of alignment at \href{https://www.overleaf.com/learn/latex/Aligning_equations_with_amsmath}{\textit{Overleaf documentation on aligning equations with amsmath}}). 

Using \verb|\label{equation}| you can refer to corresponding equation (e.g., \eref{lim}) by number.

In addition to the standard \verb"\ref{<label>}", the \tref{abrefs} provides alternative commands for quickly writing cross-references.

\begin{table}[h]
	\centering
\caption{Alternatives to the  {\tt $\backslash$ref} command for writing cross-references, as defined in the {\tt jpconf.cls} style file.} 
\label{abrefs}
\begin{tabular}{ll}
\br
Command&Result\\
\mr
\verb"\eref{<label>}"&(\verb"<num>")\\
\verb"\Eref{<label>}"&Equation~(\verb"<num>")\\
\verb"\fref{<label>}"&figure~\verb"<num>"\\
\verb"\Fref{<label>}"&Figure~\verb"<num>"\\
\verb"\sref{<label>}"&section~\verb"<num>"\\
\verb"\Sref{<label>}"&Section~\verb"<num>"\\
\verb"\tref{<label>}"&table~\verb"<num>"\\
\verb"\Tref{<label>}"&Table~\verb"<num>"\\
\br
\end{tabular}
\end{table}


\section{Figures}

Figures must be included in the source code of an article at the appropriate place in the text not grouped together at the end. 

Each figure should have a brief caption describing it and, if 
necessary, interpreting the various lines and symbols on the figure. 
As much lettering as possible should be removed from the figure itself and 
included in the caption. If a figure has parts, these should be 
labelled ($a$), ($b$), ($c$), etc. 
\Tref{blobs} gives the definitions for describing symbols and lines often
used within figure captions (more symbols are available
when using the optional packages loading the AMS extension fonts).

\begin{table}[h]
\caption{\label{blobs}Control sequences to describe lines and symbols in figure 
captions.}
\begin{center}
\begin{tabular}{lllll}
\br
Control sequence&Output&&Control sequence&Output\\
\mr
\verb"\dotted"&\dotted        &&\verb"\opencircle"&\opencircle\\
\verb"\dashed"&\dashed        &&\verb"\opentriangle"&\opentriangle\\
\verb"\broken"&\broken&&\verb"\opentriangledown"&\opentriangledown\\
\verb"\longbroken"&\longbroken&&\verb"\fullsquare"&\fullsquare\\
\verb"\chain"&\chain          &&\verb"\opensquare"&\opensquare\\
\verb"\dashddot"&\dashddot    &&\verb"\fullcircle"&\fullcircle\\
\verb"\full"&\full            &&\verb"\opendiamond"&\opendiamond\\
\br
\end{tabular}
\end{center}
\end{table}



Authors should try and use the space allocated to them as economically as possible.  
Place the figure as close as possible after the point where it is first referenced in the text. 
If there are a large number of figures it might be necessary to place some before their text citation. Figures should never appear within or after the reference list.

Individual figures should normally be centred but place two figures side-by-side if they will fit comfortably like this as it saves space. 
At times it may be convenient to put two figures side by side or the caption at the side of a figure. 
To put figures side by side, within a figure environment, put each figure and its caption into a minipage 
with an appropriate width (e.g. 3in or 18pc if the figures are of equal size) and then separate the figures 
slightly by adding some horizontal space between the two minipages (e.g. \verb"\hspace{.2in}" or \verb"\hspace{1.5pc}". 
To get the caption at the side of the figure add the small horizontal space after the \verb"\includegraphics" command and 
then put the \verb"\caption" within a minipage of the appropriate width aligned bottom, i.e. \verb"\begin{minipage}[b]{3in}" etc.

The ``\verb|figure|'' environment should be used for figures. One or
more images can be placed within a figure. 

\begin{figure}[h]
	\centering
\begin{minipage}{0.39\linewidth}
\includegraphics[width=0.95\linewidth]{img/photo_cc_.png}
\caption{\label{label5}Figure caption for first of two sided figures.}
\end{minipage}\hspace{2pc}%
\begin{minipage}{0.55\linewidth}
\includegraphics[width=0.95\linewidth]{img/name.pdf}
\caption{\label{label6}Figure caption for second of two sided figures.}
\end{minipage} 
\end{figure}

\begin{figure}[h]
\includegraphics[width=14pc]{img/name.eps}\hspace{2pc}%
\begin{minipage}[b]{14pc}\caption{\label{label7}Figure caption for a narrow figure where the caption is put at the side of the figure.}
\end{minipage}
\end{figure}

Your figures should contain a caption which describes the figure to
the reader (see \fref{fig-0}). Figure captions go below the figure. Your figures should
also include a description suitable for screen readers, to
assist the visually-challenged to better understand your work.

\begin{figure}[h]
	\centering
	\includegraphics[width=\linewidth]{img/franklinmodeld}
	\caption{Mrs. F. S. Bliven in auto (circa 1908).}
	\label{fig-0}
\end{figure}


For figures with fixed position in text use syntax of \fref{fig-0}:
\begin{verbatim}
\begin{figure}[h]
  \centering
  \includegraphics[width=0.75\linewidth]{img/franklinmodeld}
  \caption{Mrs. F. S. Bliven in auto (circa 1908).}
  \label{fig-0}
\end{figure}
\end{verbatim}

If a figure has parts these should be labelled as (a), (b), (c) etc on the actual figure. Parts should not have separate captions.
(see \fref{fig5}).
\begin{verbatim}
\begin{figure}[t]
  \begin{center}
    \begin{minipage}[b]{0.47\columnwidth}
      \includegraphics[width=1\columnwidth]{name.eps}
      \begin{center}\pt(a)\end{center}
    \end{minipage}
    \hspace{0.04\columnwidth}
    \begin{minipage}[b]{0.47\columnwidth}
      \includegraphics[width=1\columnwidth]{name.eps}
      \begin{center}\pt(b)\end{center}
    \end{minipage}
  \end{center}
  \caption{\label{fig5}A caption of figure of two parts, \pt(a) and \pt(b).}
\end{figure}
\end{verbatim}

\begin{figure}[h]
\begin{center}
\begin{minipage}[b]{0.47\columnwidth}
\includegraphics[width=1\columnwidth]{img/name.eps}
\begin{center}\pt(a)\end{center}
\end{minipage}
\hspace{0.04\columnwidth}
\begin{minipage}[b]{0.47\columnwidth}
\includegraphics[width=1\columnwidth]{img/name.eps}
\begin{center}\pt(b)\end{center}
\end{minipage}
\end{center}
\caption{\label{fig5}A caption of figure of two parts, \pt(a) and \pt(b).}
\end{figure}



\subsection{Colour illustrations}
You are free to use colour illustrations for the online version of \ees{}. 
\subsubsection{Remark}
Use over 300 dpi resolution for your figures (we prefer 600 dpi).
\paragraph{One more remark}
Don't use the lossy compressed images (e.g., JPEG). 

\section{Citations and bibliographies}

As part of the production system for \ees{}, online versions of all reference lists will, wherever possible, be linked electronically using CrossRef. \textbf{It is \textit{vitally} important for all the references to be accurate and to be carefully formatted using the guidelines below, otherwise delays may be incurred and the references may not link through CrossRef.}

Two different styles of referencing are in common use: the Harvard alphabetical system and the Vancouver numerical system.  For \ees{}, the Vancouver numerical system is preferred but authors should use the Harvard alphabetical system if they wish to do so. In the numerical system references are numbered sequentially throughout the text within square brackets, like this [2], and one number can be used to designate several references.  

\subsection{Reference lists}
A complete reference should provide the reader with enough information to locate the article concerned, whether published in print or electronic form, and should, depending on the type of reference, consist of:  

\begin{itemize}
\item name(s) and initials;
\item date published;
\item title of journal, book or other publication; 
\item titles of journal articles may also be included (optional);
\item volume number;
\item editors, if any;
\item town of publication and publisher in parentheses for {\it books};
\item the page numbers.
\end{itemize}

Up to ten authors may be given in a particular reference; where 
there are more than ten only the first should be given followed by 
`{\it et al}'. If an author is unsure of a particular journal's abbreviated title it is best to leave the title in 
full. The terms {\it loc.\ cit.\ }and {\it ibid.\ }should not be used. 

Unpublished conferences and reports should generally not be included 
in the reference list and articles in the course of publication should 
be entered only if the journal of publication is known. 

A thesis submitted for a higher degree may be included 
in the reference list if it has not been superseded by a published 
paper and is available through a library; sufficient information 
should be given for it to be traced readily.

\subsection{Formatting reference lists}
Numeric reference lists should contain the references within an unnumbered section (such as \verb"\section*{References}"). 

The use of Bib\TeX{} for the preparation and formatting of one's
references is mandatory. 

The bibliography is included in your source document with this 
command, placed just before the \verb|\end{document}| command:
\begin{verbatim}
\bibliography{bibfile}
\end{verbatim}
where ``\verb|bibfile|'' is the name, without the ``\verb|.bib|''
suffix, of the Bib\TeX{} file.

\section{Bibliographic data fields}

\subsection {References to printed journal articles}

A normal reference to a journal article is constructed as follows:

\begin{verbatim}
@article{Tkachuk_Yechkalo_Semerikov_2019, 
  author={Tkachuk, Viktoriia V. and Yechkalo, Yuliia V. 
          and Semerikov, Serhiy O.}, 
  title={The research of process of applying mobile {ICT} by university 
  students: mobile testing systems and mobile means of multimedia development}, 
  journal={Educational Dimension}, 
  year={2019}, 
  volume={53},
  issue={1}, 
  pages={125-146}, 
  url={https://doi.org/10.31812/educdim.v53i1.3839}
}
\end{verbatim}

\subsection{References to \ees{}\ articles}
Each conference proceeding published in \ees{}\ will be a separate {\it volume}; 
references should follow the style for conventional printed journals. For example:
\begin{verbatim}
@article{Striuk_2022,
  url = {https://doi.org/10.1088/1742-6596/2288/1/012012},
  year = 2022,
  month = {jun},
  publisher = {{IOP} Publishing},
  volume = {2288},
  number = {1},
  pages = {012012},
  author = {A M Striuk and S O Semerikov},
  title = {Professional competencies of future software engineers in the 
           software design: teaching techniques},
  journal = {Journal of Physics: Conference Series}

}
\end{verbatim}

%%%%%%%%%%%%%%%%%%%%%%%%%%%%%%%%%%
\subsection{References to preprints}
%For preprints there are two distinct cases:
%\renewcommand{\theenumi}{\arabic{enumi}}
%\begin{enumerate}
%\item Where the article has been published in a journal and the preprint is supplementary reference information. In this case it should be presented as:
%\medskip
\begin{verbatim}
@misc{Saptsin_Soloviev,
  url = {https://doi.org/10.48550/ARXIV.0907.1142},
  author = {Saptsin, Vladimir and Soloviev, Vladimir},
  title = {Relativistic quantum econophysics - new paradigms in complex systems modelling},
  publisher = {arXiv},
  year = {2009},
}
\end{verbatim}
%\item Where the only reference available is the preprint. In this case it should be presented as
%\medskip
%\begin{verbatim}
%@unpublished{Milson:2004,
% author = {Milson, R and Coley, A and Pravda, V and Pravdova, A},
% title = {Alignment and algebraically special tensors},
% year = {2004},
% archivePrefix = {arXiv},
% eprint = {gr-qc/0401010}, 
%}
%\end{verbatim}
%\end{enumerate}

Some more examples:
\begin{itemize}
 \item Article with preprint
\begin{verbatim}
@unpublished{Sundu:2013,
 author = {Azizi, K  and Sundu, H and S\"ung\"u, J Y and Yinelek, N},
 title = {Properties of $D_{s2}^*(2573)$ charmed-strange tensor meson},
 year = {2013},
 archivePrefix = {arXiv},
 eprint = {1307.6058}, 
 subjectarea= {hep-ph},
}
\end{verbatim}
 \item Institutional preprints or technical reports
\begin{verbatim}
@techreport{Kalitkin:1975,
 author = {Kalitkin, N. N. and Kuz'mina, L. V.},
 title = {Tables of thermodynamic functions of matter at high concentration of energy},
 type={Preprint},
 number = {35}, 
 institution = {Institute of Applied Mathematics of the USSR Academy of Sciences},
 address = {Moscow}, 
 year = {1975},
}

@techreport{Kerley2003,
 author = {Kerley, G. I.},
 title = {Equations of state for titanium and {Ti6A14V} alloy},
 type = {Report},
 number = {SAND 2003-3785},
 institution = {Sandia National Laboratories},
 address = {Albuquerque, NM},
 year = {2003}
}
\end{verbatim} 
\item Patents
\begin{verbatim}
@techreport{Rutberg2004,
 author={Rutberg, {\relax Ph} G and Safronov, A A and Shiryaev, V N},
 title={Three-phase ac plasma generator},
 type={Patent},
 number={RU 2231936},
 year={2004}
}
\end{verbatim}
\end{itemize}

%\subsection{References to electronic-only journals}
%In general article numbers are given, and no page ranges, as most electronic-only journals start each article on page 1.
%
%\begin{itemize} 
%\item For {\it New Journal of Physics} (article number may have from one to three digits)
%\numrefs{1}
%\item Fischer R 2004 Bayesian group analysis of plasma-enhanced chemical vapour deposition data {\it New. J. Phys.} {\bf 6} 25 
%\endnumrefs
%\item For SISSA journals the volume is divided into monthly issues and these form part of the article number
%
%\numrefs{2}
%\item Horowitz G T and Maldacena J 2004 The black hole final state {\it J. High Energy Phys.}  	JHEP02(2004)008
%\item Bentivegna E, Bonanno A and Reuter M 2004 Confronting the IR fixed point cosmology 	with 	high-redshift observations {\it J. Cosmol. Astropart. Phys.} JCAP01(2004)001  
%\endnumrefs
%\end{itemize} 

\subsection{References to books, conference proceedings and reports}
References to books, proceedings and reports are similar to journal references:

\begin{itemize}
 \item Complete book
\begin{verbatim}
@book{Goossens:1993,
 author = {Michel Goossens and Frank Mittelbach and Alexander Samarin},
 title = {The LaTeX Companion},
 year = {1993},
 publisher ={Addison-Wesley},
 address = {Reading, MA}
}
\end{verbatim} 
\item Book in series
\begin{verbatim}
@book{Dirac:1958,
 author = {P. A. M. Dirac},
 title = {The Principles of Quantum Mechanics},
 series = {The International Series of Monographs on Physics},
 number = {27},
 edition = {4},
 publisher = {Clarendon Press},
 address = {Oxford},
 year = {1967}
}
\end{verbatim}
\item Book chapter or some part of book
\begin{verbatim}
@inbook{Nikiforov_Novikov_Uvarov2005:ch1,
 author = {Nikiforov, A. F. and Novikov, V. G. and Uvarov, V. B.},
 title = {Quantum-Statistical Models of Hot Dense Matter},
 publisher = {Birkh\"{a}user Verlag},
 address = {Basel},
 year = {2005},
 chapter = {1},
 pages = {3--28}
}
\end{verbatim}
(You can also cite any part of book using \verb"\cite[pp 110--3]{Dirac:1958}" or \verb"\cite[chapter 4, pp 98--105]{Dirac:1958}")
\item Authored chapter 
\begin{verbatim}
@incollection{Morse:1996,
 author = {M. Morse},
 title = {Supersonic beam sources},
 booktitle = {Atomic Molecular and Optical Physics},
 editor = {F. B. Dunning and R. Hulet},
 series = {Experimental Methods in the Physical Sciences},
 volume = {29},
 publisher = {Academic},
 address = {San Diego},
 year = {1996}
}
\end{verbatim}
\item Article in conference proceedings
\begin{verbatim}
@incollection{TkachukYechkaloSemerikovKislovaHladyr_2021_0,
author="Tkachuk, Viktoriia
and Yechkalo, Yuliia
and Semerikov, Serhiy
and Kislova, Maria
and Hladyr, Yana",
editor="Bollin, Andreas
and Ermolayev, Vadim
and Mayr, Heinrich C.
and Nikitchenko, Mykola
and Spivakovsky, Aleksander
and Tkachuk, Mykola
and Yakovyna, Vitaliy
and Zholtkevych, Grygoriy",
title="{Using Mobile ICT for Online Learning During COVID-19 Lockdown}",
url={https://doi.org/10.1007/978-3-030-77592-6_3},
booktitle="Information and Communication Technologies in Education, Research, 
           and Industrial Applications",
year="2021",
publisher="Springer International Publishing",
address="Cham",
pages="46--67",
isbn="978-3-030-77592-6"
}
\end{verbatim}
or \verb|@CONFERENCE| or \verb|@inproceedings|.
\end{itemize}

\subsection{Special bibliographic data fields}
\subsubsection{Journal sections.}

Under IOP style conventions, journal names should be set in italic
type.  However, for journals with multiple lettered sections, the IOP
convention is that the journal section letter should appear in roman
type after the main journal name, \textit{e.g.}, ``\textit{J.\
	Phys.\/} A''.  Most existing \BibTeX{} styles do not make special
provision for lettered sections.  Therefore, typically, the section
letter is either included as part of the journal name
\begin{verbatim}
journal = "J. Phys. A",
volume = "38",
\end{verbatim}
or as part of the volume number
\begin{verbatim}
journal = "J. Phys.",
volume = "A38",
\end{verbatim}
in the \BibTeX{} database entry.  The \texttt{iopart-num} style
instead introduces a new optional field \verb+section+ which can be
used to specify a journal section letter.  This section letter is set in
roman type.  Moreover, if the section letter already appears in
\textit{any} of the usual locations in the database entry (at the end
of the journal name, before the volume number, or after the volume
number),
\texttt{iopart-num} will recognize it and suppress its printing.
Therefore, when you are creating the
\BibTeX{} database entry for an article in a lettered journal section,
you can still include the section letter in the \verb+journal+ or
\verb+volume+ fields, for use with other \BibTeX{} styles, 
without adversely affecting the formatting for IOP journals.  For
example, the entry for ref.~\cite{caprio2005:coherent} can be
generated with
\begin{verbatim}
journal = "J. Phys. A",
section = "A",
volume = "38",
\end{verbatim}
or
\begin{verbatim}
journal = "J. Phys.",
section = "A",
volume = "A38",
\end{verbatim}
or simply
\begin{verbatim}
journal = "J. Phys.",
section = "A",
volume = "38",
\end{verbatim}
in the \BibTeX{} database entry.  Note that section names longer than a
single letter are also supported (\textit{e.g.},
``\textit{Phys. Rev.\/} ST Accel. Beams'').

\subsubsection{Journal issue numbers.}

Journal issue numbers are not customarily included in references to
journal articles under the IOP formatting conventions.
Therefore, the
\texttt{iopart-num} style ignores the
\verb+number+ field for articles.  
However, in some periodicals (such as popular magazines or certain
journal online supplements), pagination restarts from 1 with each
issue.  For such periodicals, the issue number is an essential part of
the bibliographic information needed to identify an article.  The
\texttt{iopart-num} style therefore supports an additional field
\verb+issue+ in the \BibTeX{} database entry, which can be used to 
enforce printing of the issue number.  If a value is specified for
\verb+issue+, this value is included included parenthetically after the volume
number in the reference, as in
ref.~\cite{zamfir2005:132te-beta-enam04}.

\subsubsection{Multivolume books.}

The IOP guidelines distinguish between a volume
in a {\it series}
and a volume of a {\it  multivolume book} or {\it set}.  For a volume in
a series, the series title and volume number are given in parentheses
after the book title~\cite{ex8}.
For an individual volume in a multivolume book, the book title is
given first, followed by the volume number and volume title~\cite{ex9}.

The
\texttt{iopart-num} style supports an additional field
\verb+volumetitle+ in the \BibTeX{} database entry, which can be used to 
specify the title for an individual volume of a multivolume
book, as in refs.~\cite{ex9,bohr1998:v2}.  For example, the entry for
ref.~\cite{bohr1998:v2} is generated with
\begin{verbatim}
title = "Nuclear Structure",
volume = 2,
volumetitle = "Nuclear Deformations",
\end{verbatim}
in the \BibTeX{} database entry.  In contrast, most existing \BibTeX{} styles 
allow you to reference a volume of a multivolume
book by specifying \verb+title+ and
\verb+volume+, as in
refs.~\cite{ex7,siegbahn1965:v1}, but do not provide for inclusion of
any specific title for the individual volume.  A volume in a
series~\cite{ex8,iachello2006:liealg} is indicated in
\texttt{iopart-num}, as in most other \BibTeX{}
styles, 
by specifying
\verb+title+, \verb+volume+, and
\verb+series+ in the \BibTeX{} database entry.  

\subsubsection{E-prints, collaborations, and other data fields.}

The \texttt{iopart-num} style supports several additional data fields
(\verb+collaboration+, \verb+eid+, \verb+eprint+, 
\verb+numpages+, and \verb+url+).

\subsection{A case of non-Latin source}

When non-Latin alphabet publication cited
in an English publication,
the title of the publication (e.g., book or article) in the original language need to be both transliterated and translated in English. Other bibliographic components (including authors, publisher, address and journal name) are transliterated only \cite{IA2000}:
\begin{verbatim}
@article{IA2000,
  author ={Semerikov, S. O. and Soloviov, V. M. and Teplytskyi, I. O.},
  year=2000,
  title= {Instrumentalne zabezpechennia kursu kompiuternoho modeliuvannia
         ({I}nstrumental support of the course of computer modeling)},
  journal= {Kompiuter u shkoli i simi}, 
  issue=4,
  pages={28-31},
  url={https://lib.iitta.gov.ua/704129/}
} 
\end{verbatim}

\subsection{Best practices: export citations into a \BibTeX{} file}

A good way to make your bibliography is to exclude manual creation bibliography items whenever it possible. We strongly recommend to use the ``Cite'' (export) facilities to \BibTeX{} which available in the most of OJS installations (\fref{export}a), ACM Digital Library (\fref{export}b), Scopus (\fref{export}c), IEEE Xplore (\fref{export}d), ScienceDirect (\fref{export}e), Web of Science (\fref{export}f) etc. 

\begin{figure}[h!]
	\centering
	\begin{minipage}[h]{0.23\linewidth}
		\centering
		\includegraphics[width=1\linewidth]{img/ojs}
		 (a) \\
	\end{minipage}
	\hfill
	\begin{minipage}[h]{0.76\linewidth}
		\centering{\includegraphics[width=1\linewidth]{img/acm}} \\(b)
	\end{minipage}
	\vfill
	\begin{minipage}[h]{0.45\linewidth}
		\centering{\includegraphics[width=1\linewidth]{img/scopus}} (c) \\
	\end{minipage}
	\hfill
	\begin{minipage}[h]{0.52\linewidth}
		\centering{\includegraphics[width=1\linewidth]{img/ieee}} (d) \\
	\end{minipage}
	\vfill
	\centering
	\begin{minipage}[h]{0.35\linewidth}
		\centering
		\includegraphics[width=1\linewidth]{img/sciencedirect} (e) \\
	\end{minipage}
	\hfill
	\begin{minipage}[h]{0.5\linewidth}
		\centering{\includegraphics[width=1\linewidth]{img/wos}} (f) \\
	\end{minipage}
		\caption{Export citations into a \BibTeX{} file.}
		\label{export}
\end{figure}


\subsection{Some examples}

A paginated journal article \cite{Kiv1979K91}, an enumerated journal
article \cite{Kavetskyy2021, Striuk_2022}, a monograph (whole book) \cite{Morkun}, a
monograph/whole book in a series (see 2a in spec. document)
\cite{Harel79}, a divisible-book such as an anthology or compilation
\cite{Editor00} followed by the same example, however we only output
the series if the volume number is given \cite{Editor00a} (so series
should not be present since it has no volume number), a chapter in a
divisible book \cite{Spector90}, a chapter in a divisible book in a
series \cite{Douglass98}, a multi-volume work as book \cite{Knuth97},
an article in a proceedings (of a conference, symposium, workshop for
example) (paginated proceedings article) \cite{Andler79}, a
proceedings article with all possible elements \cite{Smith10}, an
example of an enumerated proceedings article \cite{VanGundy07}, an
informally published work \cite{Harel78}, a doctoral dissertation
\cite{Clarkson85}, a master's thesis: \cite{anisi03}, a preprint: \cite{Saptsin_Soloviev}, an online
document / world wide web resource \cite{Thornburg01, Ablamowicz07,
  Poker06}, a video game (Case 1) \cite{Obama08} and (Case 2)
\cite{Novak03} and \cite{Lee05} and (Case 3) a patent
\cite{JoeScientist001}, work accepted for publication \cite{rous08},
prolific author \cite{SaeediMEJ10} and \cite{SaeediJETC10}. Other
cites might contain `duplicate' DOI and URLs (some SIAM articles)
\cite{Kirschmer:2010:AEI:1958016.1958018}. Multi-volume works as books
\cite{MR781536} and \cite{MR781537}. A couple of citations with DOIs:
\cite{2004:ITE:1009386.1010128,Kirschmer:2010:AEI:1958016.1958018}. Online
citations: \cite{TUGInstmem, Thornburg01, R, UMassCitations}.

Refs.~\cite{ex1,ex2,ex3,ex4,ex5,ex6,ex7,ex8} are based upon example entries
from the IOP guidelines.

A lot of citations can be joined by adding the \verb|\usepackage{cite}| to document preamble: \cite{
bardeen1957:bcs,
caprio2005:coherent, 
zamfir2005:132te-beta-enam04, 
rose1957:am, 
dirac1958:qm,
siegbahn1965:v1,
bell1965:coin-lifetime,
caprio2003:diss,
doePC,
Mucklow:2007,
Tkachuk_Yechkalo_Semerikov_2019,
Zeldovich_Raizer:1967,
vonNeumann:1955,
Semerikov2021,
Trius200474,
Dixon:1997,
Morkun20144,
01_VlasenkoRovenskaChumakLovianovaAchkan,
Kunze:2003,
Sundu:2013,
Kalitkin:1975,
Kerley2003,
Rutberg2004,
Neilson:2000,
Goossens:1993,
Dirac:1958,
Nikiforov_Novikov_Uvarov2005:ch1,
Morse:1996,
TkachukYechkaloSemerikovKislovaHladyr_2021_0}.

\ack
Authors wishing to acknowledge assistance or encouragement from 
colleagues, special work by technical staff or financial support from 
organizations should do so in an unnumbered Acknowledgments section 
immediately following the last numbered section of the paper. The 
command \verb"\ack" sets the acknowledgments heading as an unnumbered
section.

\section*{ORCID iDs}
We recommend that the authors add their ORCID iDs between the acknowledgments section and the reference section. For example:
\\\\
\noindent
	S O Semerikov \url{https://orcid.org/0000-0003-0789-0272}\\
	V N Soloviev \url{https://orcid.org/0000-0002-4945-202X}\\
	A E Kiv \url{https://orcid.org/0000-0002-0991-2343}\\

The command \verb"\section*{ORCID iDs}" is used to signify the start of the ORCID iDs section:

\small\begin{verbatim}
\section*{ORCID iDs}
	S O Semerikov \url{https://orcid.org/0000-0003-0789-0272}\\
	V N Soloviev \url{https://orcid.org/0000-0002-4945-202X}\\
	A E Kiv \url{https://orcid.org/0000-0002-0991-2343}\\
\end{verbatim}\normalsize

If the paper does not have an acknowledgements section, the ORCID iDs section should follow the conclusion.



\appendix
\section{Appendices}
Technical detail that it is necessary to include, but that interrupts 
the flow of the article, may be consigned to an appendix. 
Any appendices should be included at the end of the main text of the paper, after the acknowledgments section (if any) but before the reference list.
If there are two or more appendices they will be called Appendix A, Appendix B, etc. 
Numbered equations will be in the form (A.1), (A.2), etc,
figures will appear as figure A1, figure B1, etc and tables as table A1,
table B1, etc.

The command \verb"\appendix" is used to signify the start of the
appendixes. Thereafter \verb"\section", \verb"\subsection", etc, will 
give headings appropriate for an appendix:

\small\begin{verbatim}
\appendix
\section{Appendix title 1}
\section{Appendix title 2}
\section{Appendix title 3}
\end{verbatim}\normalsize

To obtain a simple heading of 
`Appendix' use the code \verb"\section*{Appendix}". If it contains
numbered equations, figures or tables the command \verb"\appendix" should
precede it and \verb"\setcounter{section}{1}" must follow it. 

\small\begin{verbatim}
\appendix
\section*{Appendix}
\setcounter{section}{1}
\end{verbatim}\normalsize

\section*{References}


%
% BibTeX or Biber users please use (the style is already called in the class, ensure that the "woc.bst" style is in your local directory)
\bibliography{example_bibliography}
%
\end{document}

% end of file example.tex
